% p.38 — الباب الموفي بالعشرين في الافعال المعتلة — تتمة (columns 7–12)
\mspage{38}
% Layout as on p.37: six numbered columns, each with category heading, باب,
% ten rows of original forms, a rule, then ten rows of forms after تعليل.
% The running header is faded in the scan and is omitted.
% TODO: col. 12 lower block, rows 2, 9, 10: hamza vs. wāw is hard to read.
\begin{sarf}{|c|c|c|c|c|c|}
\thead{(7) مهموز العين} & \thead{(8) مهموز العين} & \thead{(9) مهموز اللام} & \thead{(10) مهموز اللام} & \thead{(11) مهموز اللام} & \thead{(12) مهموز اللام}\\
\hline
سَمِعَ & كَرُمَ & فَتَحَ & ضَرَبَ & سَمِعَ & كَرُمَ\\
\hline
يَئِسَ & لَؤُمَ & قَرَأَ & هَنَأَ & سَبِئَ & جَرُؤَ\\
يَيْئَسُ & يَلْؤُمُ & يَقْرَأُ & يَهْنِئُ & يَسْبَأُ & يَجْرُؤُ\\
يَأْسًا & لَأْمًا & قِرَاءَةً & هَنَاءً & سَبْأً & جُرْئَةً\\
يَئِيْسٌ & لَئِيْمٌ & قَارِءٌ & هَانِئٌ & سَبِيْئٌ & جَرِيْئٌ\\
\x & \x & قُرِءَ & هُنِئَ & \x & \x\\
\x & \x & يُقْرَأُ & يُهْنَأُ & \x & \x\\
\x & \x & قِرَاءَةً & هَنَاءً & \x & \x\\
\x & \x & مَقْرُوْءٌ & مَهْنُوْءٌ & \x & \x\\
اِيْئَسْ & اُلْؤُمْ & اِقْرَأْ & اِهْنِئْ & اِسْبَأْ & اُجْرُؤْ\\
لَا تَيْئَسْ & لَا تَلْؤُمْ & لَا تَقْرَأْ & لَا تَهْنِئْ & لَا تَسْبَأْ & لَا تَجْرُؤْ\\
\hline
يَئِسَ & لَؤُمَ & قَرَأَ & هَنَأَ & سَبِئَ & جَرُؤَ\\
يَيْئَسُ & يَلْؤُمُ & يَقْرَأُ & يَهْنِئُ & يَسْبَاءُ & \unclear{يَجْرُوُ}\\
\unclear{يَاْسًا} & لُوْمَةً & قِرَاءَةً & هَنَاءً & سَبْاءً & جُرْئَةً\\
يَئِيْسٌ & لَئِيْمٌ & قَارِءٌ & \unclear{هَانِءٌ} & سَبِيٌّ & جَرِيٌّ\\
\x & \x & قُرِئَ & هُنِئَ & \x & \x\\
\x & \x & يُقْرَأُ & يُهْنَأُ & \x & \x\\
\x & \x & قِرَاءَةً & هَنَاءً & \x & \x\\
\x & \x & مَقْرُوٌّ & مَهْنُوٌّ & \x & \x\\
اِيْئَسْ & اُلْؤُمْ & اِقْرَأْ & اِهْنِئْ & اِسْبَاءْ & \unclear{اُجْرُوْ}\\
لَا تَيْئَسْ & لَا تَلْؤُمْ & لَا تَقْرَأْ & لَا تَهْنِئْ & لَا تَسْبَاءْ & \unclear{لَا تَجْرُوْ}\\
\end{sarf}
